\section{Esoteric Anthropology}

\begin{quotex}
The soul is all that it knows. \flright{\textsc{Aristotle}}

\end{quotex}
Although \textbf{Rene Guenon} wrote his esoteric anthropology in the context of the Vedanta (\emph{Man and his Becoming}), there is also a deep tradition to be found in the West. By ``esoteric", of course, is meant the inner constitution of the human being, known through consciousness. Although this knowledge was part of the Hermetic stream that included him, Plato made it a key component of his oeuvre. It was then expressed in a different way by Aristotle.

Moreover, this tradition was known and developed by Christians at least through the Middle Ages. Plato's version is presumed in all the writings collected in the Philokalia. Thomas Aquinas recovered Aristotle's version. This shows the continuity between the pagan era and the Christian era in the West. Furthermore, this knowledge is not substantially different from that of the Orient, as Rene Guenon points out:

\begin{quotex}
on this and many [other conceptions], the conceptions of Aristotle are in complete agreement with those of the East. 

\end{quotex}
As the epigraph expresses, the degree of being of the soul is related to its knowledge. Hence, the soul develops as self-knowledge develops. The fault in the Western notion, according to Guenon is this:

\begin{quotex}
but this affirmation, in the case of Aristotle and his successors seems to have remained purely theoretical. It must therefore be admitted that the consequences of this idea of identification by Knowledge, as far as metaphysical realization is concerned, have continued quite unsuspected in the West, with the exception of certain strictly initiatic schools, which had no point of contact with all that usually goes by the name of `philosophy'. 

\end{quotex}
However, Christianity opened up the possibility of such a realization. In other words, beyond merely theoretical knowledge of the soul, it is a Way of salvation and liberation, something lacking in paganism (except perhaps in isolated mystery schools). This is attested to in works like the \emph{Philokalia} and \emph{Unseen Warfare}, at least to those who know how to read them esoterically.

This understanding of the soul thus has two aspects: static and dynamic. In its static aspect, it is a description of the parts of the soul. Its dynamic aspect includes the way to deeper knowledge of the soul as Aristotle described. These are analogous to the sciences of anatomy and physiognomy in medicine.

This post will deal with the static aspect. It rarely can be fully grasped without having participated in proper and direct training. Hence, to most readers, what follows will strike them as a philosophical theory. For the few, it will be recognized as fact.

\paragraph{Centers}
There are three basic centers which correspond in man to the three kingdoms: the sensual, emotional, and intellectual centers. Since the soul is not in a particular place, these cannot be resolved to specific biological functions. This is the condition of man as he is. As such, they are disordered and do not serve as a true guide to life. Hence, it is necessary to transcend these three centers.

The emotional and intellectual centers have corresponding higher centers, that, for the most part, we are unaware of. The higher intellectual center is not disordered at all and is perfect. This is the ``image of God" that was left intact after the Fall. The sexual center, or eros, often mistakenly merged with the sensual center, is necessary for the second birth. Obviously, it can be part of a sensual life, but when sublimated, it can lead to transcendence. Hence, eros needs to be directed toward the logos. These centers are summarized in Table \ref{tab:CentersMan}.

\begin{table}[h]\small
\begin{tabularx}{\textwidth}{llXX}\toprule
Level &
Center &
Function &
Other terms\\\midrule
1 &
Sensual and Sexual &
This center regulates the instinctive and motor functions, and the life of sensations &
Epithumia and eros

Vegetative soul\\\midrule
2 &
Emotional &
The center corresponds to ordinary emotional life &
Thumos, Philonikon

Animal soul\\\midrule
3 &
Intellectual &
The lower intellectual center corresponds to the life of thought, i.e., of forms &
Dianoia

Ratio

Intellectual Soul\\\midrule
4 &
Higher Emotional &
This corresponds to higher emotions related to honour, loyalty, courage, and nobility. &
Philotimon\\\midrule
5 &
Higher Intellectual &
This corresponds to direct intuition, beyond words and thought. &
Nous

Intellectus

Buddhi\\\bottomrule
\end{tabularx}
\caption{Centers of Man}
\label{tab:CentersMan}
\end{table}
This outline leaves out important details. The three lower centers each have a positive and a negative aspect. For example, the sensual center may be known through automatic, or instinctive, attraction and avoidance. Everyone experiences positive and negative emotions. The intellectual center will agree or disagree with an idea.

Furthermore, the three centers or souls interpenetrate each other. Hence, there will be certain emotions associated with sensual activity, or instinctive functions will display intelligence. All that is too complex for this introduction.

\paragraph{Exterior Men}
Everyone is dominated naturally by one of the three lower centers; this is fixed and cannot be changed. At this point of non-development, such are called ``exterior men or women" because their attention is dominated by external facts, conditions, or events. Julius Evola refers to this initial state in these terms:

\begin{quotex}
From its first moment, we can say that the ``I" still lives only as if in a dream: it is not yet a self-consciousness, nor an autonomous principle of action: immersed in an immediate, indistinct coalescence with nature and the world, we can say that it is not so much he who thinks, speaks, and asserts himself, but rather that various forces and impulses think, speak and assert themselves in him. He therefore is only a type of medium, a passive instrument that has his very life outside himself, and he experiences everything as grace, as spontaneity, as the immediate self-manifestation of something that transcends him. 

\end{quotex}
At this stage, the problem of certainty does not exist, since whatever arises spontaneously in consciousness is presumed to be a true image of the world. That is why initiation must begin by provisionally questioning all one's cherished beliefs and opinions. Unfortunately, apart from those very few born righteous, no one is willing to do that without first having undergone some sort of spiritual, emotional, or intellectual crisis. The three centers give rise to different misconceptions.

\begin{itemize}
\item \textbf{Sensual misconceptions}. Those dominated by the first or sensual center regard the purpose of life as the pursuit and satisfaction of sensual desires. 
\item \textbf{Emotional misconceptions}. Those dominated by the second or emotional center consider their emotional life as guide to life. Hence, for example, they may be overly concerned about taking offence. If they hold a particular opinion, say a political or religious one, they believe that the degree of passion associated with that opinion is an indicator of its validity. 
\item \textbf{Intellectual misconceptions}. Those dominated by the third or intellectual center are attracted to ideas, ideologies, or opinions for their own sake. For them, questions of epistemology, i.e., ``how do I know what I claim?", don't enter their consciousness. They are prone to frequent changes or conversions. They are recognized by the tendency to categorize, and even dismiss, the giants of thought with little more than a one sentence summary. The intent is to sound sophisticated in their peer group. The danger of this type for Tradition is that they will try to understand Tradition philosophically (as Guenon pointed out, see above), or they will misunderstand esoteric texts by reducing them to their limited frame of reference. 
\end{itemize}
Of course, this just scratches the surface of possible misconceptions. Those who actually follow an esoteric way will learn to recognize these misconceptions both in themselves and in others.

\paragraph{Interior Men}
The few who manage to escape the attachment to exteriority will be obsessed with finding a way out. Julius Evola describes the beginning of this phase:

\begin{quotex}
Now in this second stage, following the ``period of spontaneity", the ``I" turns toward autonomy and individual existence – that consequently happens only in so far as the original connection with the whole is gradually broken, and as a gradual tearing away from the World takes place. Consequently, what used to be familiar to the individual is made alien and impenetrable, what intuitive certainty used to reveal to him as indisputable fact is made dubious and problematic. 

\end{quotex}
So the first task is to begin observing and transcending the misconceptions of the first stage. A temporary ``I" develops that integrates the three lower centers. This is still a transitional, or psychological, state, often represented as the lower Ego or the Self in Jungian terms.

Beyond that, there are three further stages: the development of the Real I, Consciousness, and True Will. This are all the possibilities available to man as such and depend on recentering on the higher emotional and intellectual centers. The final stage depends on the complete sublimation of the sexual center. But this process will have to await the discussion of the dynamic aspect of esoteric anthropology.



\flrightit{Posted on 2014-11-09 by Cologero }

\begin{center}* * *\end{center}

\begin{footnotesize}\begin{sffamily}



\texttt{TeenyTinyTarot on 2014-11-10 at 08:30 said: }

\url{http://nazbol.net/library/authors/Rene\%20Guenon/Rene\_Guenon\_Man\_and\_His\_Becoming\_according\_to\_the\_Vedanta\_Collected\_Works\_of\_Rene\_Guenon\_\_2004.pdf}


\hfill

\texttt{Paulo on 2014-11-10 at 11:29 said: }

``From its first moment, we can say that the ``I" still lives only as if in a dream: it is not yet a self-consciousness, nor an autonomous principle of action: immersed in an immediate, indistinct coalescence with nature and the world, we can say that it is not so much he who thinks, speaks, and asserts himself, but rather that various forces and impulses think, speak and assert themselves in him. He therefore is only a type of medium, a passive instrument that has his very life outside himself, and he experiences everything as grace, as spontaneity, as the immediate self-manifestation of something that transcends him."

Should we see any difference between Jungian process of individuation,Gurdjieff fourth way,Ramana Maharishi,Nisargadattha Maharaj self enquiry and the traditionalist way?


\hfill

\texttt{Avery on 2014-11-13 at 11:23 said: }

``it is a Way of salvation and liberation, something lacking in paganism (except perhaps in isolated mystery schools)"

I realize that I am making a rather foolish complaint given the developed stage of this blog, but this is the trouble that gave rise to Guenonian Traditionalism in the first place. Is the Truth available to pagans only through secret rituals? In other words, is the Christian God a Gnostic God for those without a covenant? I don't think this is a minor point, since it would appear to upset the objective of catholic (universal) religion.

If you've written about this elsewhere, feel free to link me to there.


\hfill

\texttt{Cologero on 2014-11-13 at 13:22 said: }

I don't believe I am grasping the issue, Avery. There is no ``Christian God", just the God; if I ever gave another impression, I'll have to clarify or retract it.

I don't ``know" about the pagans directly, just what can be inferred through their writings. Clearly, pagans from Hermes to Plato and beyond knew about the God and even the Logos.

The populace, on the other hand, was taught polytheism with gods selectively interfering in human affairs. There was no solution to the problem of death; descendants were expected to keep the flame going. I had in mind some passages from Bachofen but the book is not at hand right now.

What do you suppose was the purpose of those mystery schools (and why did you change the designation to ``secret rituals")? Why did the early Christians believe that they superseded those schools since they were the esoteric teaching.

What ``gave rise to Guenonian Traditionalism" was Guenon's failure to make contact with an esoteric tradition in the West. Isn't that because the original distinctions between the catachumens, the faithful, and the perfect became blurred? And aren't most exoteric christians, including academics, more like pagans except they believe in one god rather than many?

I have quoted on several occasions (feel free to find the links yourself) those theologians who regarded previous traditions as precursors, not as indicating a rupture.

Given this, please be specific about your complaint so we can work it out.


\end{sffamily}\end{footnotesize}
